\documentclass[11pt]{article}
\usepackage[margin=1in]{geometry}
\usepackage{amsmath,amssymb,amsthm}
\usepackage{booktabs}
\usepackage{enumitem}
\usepackage{hyperref}
\usepackage{graphicx}
\usepackage{xcolor}

% ---------------------------------------------------------------
% Paper/note toggle: compile with \papertrue for paper-ready output
% (hides development notes, intro, §7 connection, verbose phases)
% Default: \papertrue = full working note
% ---------------------------------------------------------------
\newif\ifpaper
\papertrue   % change to \papertrue to get paper-ready section

\newtheorem{theorem}{Theorem}
\newtheorem{proposition}[theorem]{Proposition}
\newtheorem{corollary}[theorem]{Corollary}
\newtheorem{lemma}[theorem]{Lemma}
\newtheorem{definition}[theorem]{Definition}
\theoremstyle{remark}
\newtheorem{remark}[theorem]{Remark}

\ifpaper
  \title{Adaptive Bandwidth, Score Homogeneity, and\\
    Conditional Coverage in CKME-CP}
  \author{}
  \date{}
\else
  \title{Note~2: Score Homogeneity Mechanism\\
    {\large Adaptive Bandwidth, Uniform Bias, and Conditional Coverage}}
  \author{Working Note --- CKME-CP Project}
  \date{April 2026}
\fi

\begin{document}
\maketitle

\begin{abstract}
We identify a new regime for conformal prediction with CDF-based
scores: when the bandwidth of the smooth indicator function is
proportional to the local noise scale, the conformal score distribution
becomes approximately $x$-independent, even though the underlying CDF
estimate is biased. This ``score homogeneity'' property ensures that a
single calibration threshold provides uniform conditional coverage ---
a guarantee unavailable to fixed-bandwidth estimators or
quantile-first methods. We formalize the mechanism through the concept
of \emph{effective bandwidth in standardized scale}, prove an
over-smoothed collapse proposition showing what fails when homogeneity
is achieved trivially, and validate the theory with a three-phase
experiment covering mild, extreme, and non-Gaussian heteroscedasticity.
\end{abstract}

\tableofcontents
\newpage


%% ====================================================================
% NOTE-ONLY: introduction is replaced by a brief paragraph in paper mode
\ifpaper\else
\section{Introduction}
\label{sec:intro}
%% ====================================================================

Split conformal prediction (CP) calibrates a prediction set by
choosing a threshold $\hat{q}$ such that
$P(s(X, Y) \le \hat{q}) \ge 1 - \alpha$ on fresh calibration data.
This guarantees \emph{marginal} coverage, but conditional coverage
$P(Y \in C(x) \mid X = x) \ge 1 - \alpha$ requires additional
structure.

Existing theory identifies two regimes:
\begin{enumerate}
  \item \textbf{Regime~1 (consistency):} If the score converges to the
    oracle score uniformly, conditional coverage follows. This requires
    the estimator to be consistent.
  \item \textbf{Regime~2 (no structure):} Without consistency or
    additional structure, only marginal coverage holds.
\end{enumerate}

We identify a \textbf{third regime}: the CDF estimator may be biased,
but if the bias structure is \emph{spatially homogeneous} --- i.e., the
score distribution $F_{s|x}$ does not depend on $x$ --- then the single
CP threshold $\hat{q}$ simultaneously satisfies the conditional
coverage inequality at every $x$. We call this \textbf{score
homogeneity}.

The key insight is that score homogeneity can be achieved by a simple
design choice: setting the CKME bandwidth proportional to the local
noise scale, $h(x) = c\,\sigma(x)$. This makes the effective bandwidth
constant in standardized coordinates, producing a common smoothed CDF
estimate for all $x$.

This note formalizes the mechanism, proves when it succeeds and when it
degenerates, and validates the theory experimentally.
\fi % end \ifpaper\else for Introduction


%% ====================================================================
\section{Setup and Notation}
\label{sec:notation}
%% ====================================================================

\paragraph{Data model.}
$Y = f(x) + \sigma(x)\varepsilon$, where $f(x) = E[Y \mid X = x]$,
$\sigma(x) > 0$ is the conditional noise scale, and
$\varepsilon \sim G(\cdot)$ with $E[\varepsilon] = 0$,
$\mathrm{Var}(\varepsilon) = 1$, independent of $x$.

\paragraph{CKME CDF estimator.}
Given training data $(X_i, Y_{ij})$ ($i = 1,\ldots,n$;
$j = 1,\ldots,r$), CKME estimates:
\begin{equation}
\label{eq:ckme-cdf}
  \hat{F}(y \mid x) = \sum_{i=1}^{n} c_i(x) \left[
    \frac{1}{r} \sum_{j=1}^{r} \Psi\!\left(\frac{y - Y_{ij}}{h}\right)
  \right],
\end{equation}
where $\Psi(u) = 1/(1+e^{-u})$ is the logistic CDF,
$h > 0$ is the bandwidth, and $c_i(x)$ are kernel regression weights
satisfying $\sum_i c_i(x) \approx 1$.

\paragraph{Conformal score.}
The abs-median score is:
\begin{equation}
\label{eq:score}
  s(x, y) = |\hat{F}(y \mid x) - 0.5|.
\end{equation}

\paragraph{Standardized coordinates.}
Define $z = (y - f(x))/\sigma(x)$ and
$Z_{ij} = (Y_{ij} - f(x))/\sigma(x)$.


%% ====================================================================
\section{Adaptive Bandwidth and Score Homogeneity}
\label{sec:mechanism}
%% ====================================================================

% ------------------------------------------------------------------
\subsection{The Effective Bandwidth}
\label{sec:h-eff}
% ------------------------------------------------------------------

The smooth indicator in~\eqref{eq:ckme-cdf} evaluates:
\[
  \Psi\!\left(\frac{y - Y_{ij}}{h}\right)
  = \Psi\!\left(\frac{(y - f(x)) - (Y_{ij} - f(x))}{h}\right)
  = \Psi\!\left(\frac{\sigma(x)(z - Z_{ij})}{h}\right)
  = \Psi\!\left(\frac{z - Z_{ij}}{h/\sigma(x)}\right).
\]

\begin{definition}[Effective bandwidth]
\label{def:h-eff}
The \textbf{effective bandwidth in standardized scale} is
\begin{equation}
\label{eq:h-eff}
  h_{\mathrm{eff}}(x) \;=\; \frac{h}{\sigma(x)}.
\end{equation}
It controls the smoothing applied to the standardized noise CDF
$G(z)$: a larger $h_{\mathrm{eff}}$ produces a smoother (more biased)
estimate.
\end{definition}

\paragraph{Intuition.}
The bandwidth $h$ operates in the original $y$-scale, but the ``natural
unit'' for smoothing is the noise scale $\sigma(x)$. When $h$ is
fixed, the estimator over-smooths where noise is small and
under-smooths where noise is large.

% ------------------------------------------------------------------
\subsection{Fixed vs.\ Adaptive Bandwidth}
\label{sec:fixed-vs-adaptive}
% ------------------------------------------------------------------

\begin{proposition}[Score homogeneity under adaptive bandwidth]
\label{prop:homogeneity}
Under the location-scale model $Y = f(x) + \sigma(x)\varepsilon$
with $\varepsilon \sim G(\cdot)$ independent of $x$:
\begin{enumerate}[label=(\alph*)]
  \item \textbf{Fixed $h$:}
    $h_{\mathrm{eff}}(x) = h/\sigma(x)$ varies with $x$.
    The CDF estimate in standardized scale is
    $\hat{G}_{h/\sigma(x)}(z)$, which depends on $x$ through
    $\sigma(x)$. The score distribution $F_{s|x}$ is $x$-dependent.
  \item \textbf{Adaptive $h(x) = c\,\sigma(x)$:}
    $h_{\mathrm{eff}}(x) = c$ (constant). The CDF estimate becomes
    \begin{equation}
    \label{eq:G-hat-c}
      \hat{F}(y \mid x) = \sum_i c_i(x)\left[
        \frac{1}{r}\sum_j \Psi\!\left(\frac{z - Z_{ij}}{c}\right)
      \right] \approx \hat{G}_c(z),
    \end{equation}
    a common kernel-smoothed estimate of $G(z)$ with bandwidth $c$,
    independent of $x$. Consequently, the score
    $s(x,Y) = |\hat{G}_c(Z) - 0.5|$ has an $x$-independent distribution.
\end{enumerate}
\end{proposition}

\paragraph{Intuition.}
Adaptive bandwidth ``standardizes'' the smoothing: every $x$ sees the
same effective resolution $c$ when measured in units of the local noise.
The CDF bias $\hat{G}_c(z) - G(z)$ is a function of $c$ and $G$ only,
not of $x$.

\begin{proof}[Proof sketch]
Under $h(x) = c\,\sigma(x)$, the argument of $\Psi$ becomes
$(z - Z_{ij})/c$ as shown above. Since $Z_{ij} \sim G(\cdot)$
independently of $x$, the site-averaged indicator
$\frac{1}{r}\sum_j \Psi((z - Z_{ij})/c)$ is a kernel density estimate
of $G$ at $z$ with bandwidth $c$, using i.i.d.\ samples from $G$.
The dependence on $x$ enters only through $c_i(x)$, which are kernel
regression weights satisfying $\sum_i c_i(x) \approx 1$. Under
standard kernel regression consistency conditions,
$\hat{F}(y|x) \to \hat{G}_c(z)$ as $n \to \infty$, where
$\hat{G}_c$ is the pooled kernel-smoothed CDF of $G$.
\end{proof}

\begin{remark}[Why CDF-first is essential]
\label{rem:cdf-first}
Quantile-first scores such as
$s(x,y) = \max(\hat{q}_{\mathrm{lo}}(x) - y, y - \hat{q}_{\mathrm{hi}}(x))$
do not have PIT structure. Even if the quantile estimator has
$x$-homogeneous bias, the score distribution depends on the local
noise scale through the $y$-space comparison. CDF-first scores, by
operating in probability space $[0,1]$, automatically inherit the
standardization that makes homogeneity possible.
\end{remark}


% ------------------------------------------------------------------
\subsection{Over-Smoothed Collapse}
\label{sec:oversmooth}
% ------------------------------------------------------------------

Score homogeneity alone is insufficient: trivially setting $h \to \infty$
also makes $\hat{F}(y|x)$ independent of $x$ (all CDFs flatten to
$1/2$), but the resulting CP intervals are useless.

\begin{proposition}[Over-smoothed collapse]
\label{prop:oversmooth}
In the regime $h \to \infty$ (equivalently, fixed $h$ with
$\sigma(x) \ll h$ for all $x$), the abs-median CP interval collapses
to a global constant width:
\begin{equation}
\label{eq:collapsed-width}
  W(x) = U(x) - L(x) \;\to\;
  2\,Q_{1-\alpha}\bigl(\{|Y_j - \hat{f}(X_j)|\}_{j=1}^{n_{\mathrm{cal}}}\bigr)
  \quad \text{for all } x,
\end{equation}
where $Q_{1-\alpha}$ denotes the $(1-\alpha)$-quantile of the
calibration residuals.
\end{proposition}

\begin{proof}[Proof sketch]
The logistic CDF admits the Taylor expansion
$\Psi(u) \approx \frac{1}{2} + \frac{u}{4}$ for $|u| \ll 1$.
When $h$ is large, $u = (y - Y_{ij})/h$ is small, so:
\[
  \hat{F}(y \mid x)
  \approx \sum_i c_i(x)\!\left[\frac{1}{2}
    + \frac{y - \bar{Y}_i}{4h}\right]
  = \frac{1}{2} + \frac{y - \hat{f}(x)}{4h},
\]
where $\hat{f}(x) = \sum_i c_i(x)\bar{Y}_i$ is the kernel regression
estimate. The abs-median score becomes:
\[
  s(x, y) = \left|\hat{F}(y|x) - \frac{1}{2}\right|
  = \frac{|y - \hat{f}(x)|}{4h}.
\]
The CP threshold $\hat{q}$ is the $(1-\alpha)$-quantile of
$\{|Y_j - \hat{f}(X_j)|/(4h)\}$. The prediction interval
$\{y : s(x,y) \le \hat{q}\}$ requires
$|y - \hat{f}(x)|/(4h) \le \hat{q}$, giving:
\[
  L(x) = \hat{f}(x) - 4h\hat{q}, \qquad
  U(x) = \hat{f}(x) + 4h\hat{q}.
\]
The width $U(x) - L(x) = 8h\hat{q}$. But
$\hat{q} = Q_{1-\alpha}(|Y_j - \hat{f}(X_j)|)/(4h)$, so
$8h\hat{q} = 2\,Q_{1-\alpha}(|Y_j - \hat{f}(X_j)|)$.
The factor $h$ cancels exactly, and the width depends only on the
marginal distribution of calibration residuals.
\end{proof}

\begin{remark}[Width floor and coverage distortion]
\label{rem:width-distortion}
The collapsed width
$W^* = 2\,Q_{1-\alpha}(|\sigma(X)\varepsilon|)$ depends on the
\emph{marginal} distribution of $\sigma(X)$, not the local
$\sigma(x)$:
\begin{itemize}
  \item Where $\sigma(x) \ll E[\sigma(X)]$: the interval is vastly
    too wide $\to$ severe over-coverage (wasted precision).
  \item Where $\sigma(x) \gg E[\sigma(X)]$: the interval is too narrow
    $\to$ under-coverage.
\end{itemize}
\end{remark}

\begin{remark}[Numerical example: \texttt{exp2}]
\label{rem:exp2-collapse}
For \texttt{exp2}: $\sigma(x) = \sqrt{0.01 + 0.2(x-\pi)^2}$,
$x \in [0, 2\pi]$, $\alpha = 0.1$.
\begin{itemize}
  \item At $x = \pi$: $\sigma(\pi) = 0.1$, oracle 90\% interval width
    $\approx 2 \times 1.645 \times 0.1 = 0.33$.
  \item At $x = 0$: $\sigma(0) \approx 2.0$, oracle width
    $\approx 6.53$.
  \item Collapsed width:
    $W^* = 2\,Q_{0.9}(|\sigma(X)\varepsilon|) \approx 2.95$ (computed
    from the marginal residual distribution with $X \sim \mathrm{Uniform}[0,2\pi]$).
  \item At $x = \pi$: $W^*/W_{\mathrm{oracle}} \approx 9.0$
    ($9\times$ too wide).
  \item At $x = 0$: $W^*/W_{\mathrm{oracle}} \approx 0.45$
    ($55\%$ under-coverage risk).
\end{itemize}
\end{remark}

% ------------------------------------------------------------------
\subsection{The Two Conditions for Useful Homogeneity}
\label{sec:two-conditions}
% ------------------------------------------------------------------

The over-smoothed collapse shows that score homogeneity is
\emph{necessary but not sufficient}. Two conditions are needed
simultaneously:

\begin{enumerate}
  \item \textbf{Score homogeneity:}
    $d_{\mathrm{TV}}(F_{s|x}, F_{s|x'}) \approx 0$ for all $x, x'$.
    This ensures a single $\hat{q}$ works uniformly.
  \item \textbf{Local adaptivity:}
    The interval width $W(x)$ reflects local $\sigma(x)$.
    This ensures the interval is not trivially constant.
\end{enumerate}

Adaptive bandwidth $h(x) = c\,\sigma(x)$ with moderate $c$ achieves
both: the CDF is smoothed uniformly in standardized scale (condition~1),
but the smoothing preserves the $\sigma(x)$-dependence of the
quantiles because the CDF estimate still distinguishes different
$\sigma(x)$ through the $y$-to-$z$ transformation (condition~2).

Fixed $h$ with $h \to \infty$ achieves condition~1 only, violating
condition~2 (Proposition~\ref{prop:oversmooth}).
Fixed $h$ with $h \to 0$ achieves condition~2 only (sharp CDF,
locally adaptive width), but the score distribution is
$x$-dependent, violating condition~1.

% ------------------------------------------------------------------
\subsection{Connection to CP$^2$ Framework}
\label{sec:cp2}
% ------------------------------------------------------------------

Our score homogeneity result connects to the conditional conformal
prediction literature through the total variation framework
\cite{Barber2023}. Under approximate exchangeability, the conditional
coverage satisfies
\begin{equation}
\label{eq:cp2-bound}
  P(Y \in C(x) \mid X = x) \ge 1 - \alpha
  - d_{\mathrm{TV}}(x) - p_{n+1}(x),
\end{equation}
where $d_{\mathrm{TV}}(x)$ measures the discrepancy between the
conditional and marginal score distributions and $p_{n+1}(x)$ is an
exchangeability correction. The CP$^2$ framework \cite{CP2} provides
RKHS-based bounds on a similar coverage gap.

\paragraph{When $d_{\mathrm{TV}}(x) = \delta$ (constant).}
The bound becomes $P(Y \in C(x) \mid X=x) \ge 1 - \alpha - \delta
- p_{n+1}$, a \emph{uniform} bound independent of $x$.

\paragraph{Our contribution.}
CP$^2$ does not identify when $d_{\mathrm{TV}}(x)$ is $x$-independent.
We provide a concrete sufficient condition:

\begin{lemma}[Adaptive $h$ implies uniform $d_{\mathrm{TV}}$]
\label{lem:uniform-dtv}
Under the location-scale model with
$h(x) = c\,\sigma(x)$, the CDF-based conformal score distribution
satisfies $d_{\mathrm{TV}}(F_{s|x}, F_{s|x'}) \to 0$ for all
$x, x'$ as $n \to \infty$. Hence $d_{\mathrm{TV}}(x)$ (relative to
the pooled score distribution) is asymptotically constant in $x$.
\end{lemma}

\paragraph{Important subtlety.}
Uniform $d_{\mathrm{TV}}$ alone is not sufficient (Section~\ref{sec:oversmooth}).
The combination of small constant $d_{\mathrm{TV}}$ \emph{and} locally
adaptive intervals requires the bandwidth to be proportional to
$\sigma(x)$, not merely large.

% ------------------------------------------------------------------
% Promoted to top-level section so it reads as the main theorem section
\section{Coverage Bound and Rate Analysis}
\label{sec:error-comparison}
% ------------------------------------------------------------------

RCP (Rectified Conformal Prediction) \cite{RCP} and our CKME-based
approach both aim to reduce the $x$-dependence of conformal score
distributions, but they operate on fundamentally different error
decompositions. We make this distinction precise.

\paragraph{RCP's error structure.}
RCP's strategy is to transform the conformal score via a rectification
map $s \mapsto \tilde{s} = T(s, x)$ so that the transformed score's
conditional $\alpha$-quantile $\tau(x) :=
Q_\alpha(\tilde{s} \mid X = x)$ becomes $x$-independent.
In practice, $\tau(x)$ is unknown and must be estimated. Let
$\hat\tau(x)$ denote the estimated conditional quantile (obtained,
e.g., via local quantile regression with kernel bandwidth $b$).
RCP's conditional coverage theorem \cite{RCP} shows that the
coverage gap is controlled by
\begin{equation}
\label{eq:rcp-error}
  \bigl|P(Y \in C(x) \mid X=x) - (1-\alpha)\bigr|
  \;\lesssim\; \sup_x |\hat\tau(x) - \tau(x)|,
\end{equation}
i.e., the \emph{quantile-of-score estimation error}.
This error in turn inherits the usual bias--variance trade-off of
the local quantile regression bandwidth $b$: small $b$ yields low bias
but high variance; large $b$ reverses the trade-off. The bandwidth $b$
is a nuisance parameter that must be tuned, and its optimal value
depends on the smoothness of $\tau(x)$.

\paragraph{Our error structure.}
In the CKME framework with CDF-based scores, we do not estimate a
conditional quantile of the score. Instead, the relevant error
decomposition is:
\begin{equation}
\label{eq:our-error}
  d_{\mathrm{TV}}\bigl(F_{s|x},\; F_{s,\mathrm{pivot}}\bigr)
  \;\le\;
  \underbrace{\frac{\pi^2 c^2}{6}\,\|g''\|_1}_{\text{(i) smoothing bias}}
  \;+\;
  \underbrace{\sup_y\,|\hat F(y\mid x) - F(y\mid x)|}_{\text{(ii) CDF estimation error}},
\end{equation}
where term~(i) arises from the logistic smooth indicator with adaptive
bandwidth $h(x) = c\,\sigma(x)$ (Devroye and Gy\"{o}rfi \cite{DG85}),
and term~(ii) is the CKME regression error (which is $O_p(n^{-s})$
for some rate $s > 0$ depending on kernel smoothness and dimension).

When $\sigma(x)$ is unknown and replaced by $\hat\sigma(x)$, a third
term appears:
\begin{equation}
\label{eq:plugin-error}
  \text{(iii) plug-in error:}\quad
  |\hat h(x) - h(x)| = c\,|\hat\sigma(x) - \sigma(x)|,
\end{equation}
which propagates into both terms~(i) and~(ii) through the
bandwidth. Bounding this term quantitatively is addressed in
Section~\ref{sec:plugin-rate}.

\paragraph{Structural comparison.}
The key difference is summarized as follows:

\begin{center}
\small
\begin{tabular}{@{}lll@{}}
\toprule
  & \textbf{RCP} & \textbf{CKME (ours)} \\
\midrule
Core estimated object
  & $\hat\tau(x)$: conditional quantile of score
  & $\hat F(y \mid x)$: conditional CDF \\
Coverage-controlling error
  & $\sup_x |\hat\tau(x) - \tau(x)|$
  & $d_{\mathrm{TV}}(F_{s|x}, F_{s,\mathrm{pivot}})$ \\
Bandwidth role
  & nuisance parameter in local QR
  & design parameter with explicit bias formula \\
Error decomposition
  & QR bias--variance (implicit in $b$)
  & smoothing bias $O(c^2)$ + CDF error $O_p(n^{-s})$ \\
Plug-in parameter
  & $\tau(x)$ itself
  & $\sigma(x)$ (noise scale) \\
\bottomrule
\end{tabular}
\end{center}

\begin{remark}[Different ``quantile errors'']
\label{rem:quantile-errors}
Both frameworks involve quantile estimation, but the quantile objects
are distinct. RCP estimates $Q_\alpha(\tilde{s} \mid X=x)$, the
$\alpha$-quantile of the \emph{transformed score} conditional on $x$.
CKME, when used for prediction intervals, estimates
$Q_\tau(Y \mid X=x)$, the $\tau$-quantile of the \emph{response}
conditional on $x$, obtained by inverting $\hat F(y \mid x)$.
The former is an intermediate quantity for calibration; the latter is
the final inferential target. This distinction is important: our
quantile error enters through CDF inversion and is bounded by the
CDF estimation error via $|Q_\tau - \hat Q_\tau| \le \|\hat F - F\|_\infty / f(Q_\tau \mid x)$,
whereas RCP's quantile error is a direct estimation problem in its own
right.
\end{remark}

\paragraph{Implications for bandwidth selection.}
In RCP, the bandwidth $b$ of the local quantile regression is a
nuisance parameter: it controls the estimation quality of $\hat\tau(x)$
but has no structural role in the score distribution itself.
In our framework, the bandwidth $c$ of the smooth indicator plays a
dual role: it is both an estimation parameter (controlling CDF
smoothing bias) and a \emph{design} parameter (controlling score
homogeneity). The explicit decomposition~\eqref{eq:our-error}
makes this dual role transparent and provides a principled basis for
selecting $c$, which has no direct analogue in the RCP framework.

\paragraph{Two-sided conditional coverage bound.}
A structural advantage of the $d_{\mathrm{TV}}$ framework is that it
yields a \emph{two-sided} bound on conditional coverage. Since total
variation controls discrepancy in both directions, we obtain:

\begin{theorem}[Two-sided conditional coverage bound]
\label{prop:two-sided}% label kept for back-compat; cite as Theorem~\ref{prop:two-sided}
Let $\hat{q}$ be the split-conformal calibration quantile computed
from $n$ calibration scores. Suppose
$d_{\mathrm{TV}}(F_{s|x}, F_{s,\mathrm{pool}}) \le \delta(c, n)$
for all $x$, where $F_{s,\mathrm{pool}}$ is the marginal (pooled)
score distribution. Then
\begin{equation}
\label{eq:two-sided}
  \bigl|P(Y \in C(x) \mid X = x) - (1 - \alpha)\bigr|
  \;\le\; \delta(c, n) + p_{n+1},
\end{equation}
where $p_{n+1} = O(1/n)$ is the finite-sample calibration correction.
Under the adaptive bandwidth $h(x) = c\,\sigma(x)$ with the
location-scale model, the bound~\eqref{eq:our-error} gives
\[
  \delta(c, n) \;\le\;
  \frac{\pi^2 c^2}{6}\,\|g''\|_1 + O_p(n^{-s}).
\]
\end{theorem}

\begin{proof}[Proof sketch]
By definition of $d_{\mathrm{TV}}$,
$|P(s \le \hat{q} \mid X=x) - P(s \le \hat{q})| \le
d_{\mathrm{TV}}(F_{s|x}, F_{s,\mathrm{pool}})$.
Standard split-conformal calibration gives
$|P(s \le \hat{q}) - (1 - \alpha)| \le p_{n+1}$.
The triangle inequality yields~\eqref{eq:two-sided}.
The bound on $\delta(c,n)$ follows from~\eqref{eq:our-error}.
\end{proof}

\begin{remark}[Two-sided vs.\ one-sided bounds]
\label{rem:two-sided}
RCP's coverage theorem is essentially one-sided: it guarantees that
conditional coverage does not fall too far below $1 - \alpha$
(a safety guarantee), but does not directly control over-coverage.
Proposition~\ref{prop:two-sided} controls both directions:
\begin{itemize}
  \item The \textbf{lower bound} ensures safety (no severe
    under-coverage at any $x$).
  \item The \textbf{upper bound} ensures efficiency (no severe
    over-coverage, which would imply unnecessarily wide intervals
    at some $x$).
\end{itemize}
Together, they assert that conditional coverage is approximately
$1 - \alpha$ \emph{uniformly} over $x$ --- a near-conditional-validity
statement. This two-sided nature is a direct consequence of using
$d_{\mathrm{TV}}$ (an absolute discrepancy) rather than a one-sided
quantile perturbation argument.

The absolute-value form~\eqref{eq:two-sided} also corresponds
directly to the experimental metric \texttt{cov\_gap\_sup}
$= \sup_x |\text{coverage}(x) - (1-\alpha)|$ reported in the
empirical results, making the theory--experiment connection
transparent.
\end{remark}

\subsection{Plug-in Bandwidth for the Noise Scale}
\label{sec:plugin-rate}

The theoretical adaptive rule $h(x) = c\,\sigma(x)$ requires the
unknown noise scale $\sigma(x)$. In practice we replace it by a
plug-in estimator $\hat\sigma(x)$, which introduces a third error
component, $|\hat h(x) - h(x)| = c\,|\hat\sigma(x)-\sigma(x)|$, that
was left unquantified in~\eqref{eq:plugin-error}. Here we derive
its rate when $\hat\sigma$ is built from replicated Stage-1 data.

\paragraph{Two-step construction of $\hat\sigma$.}
Assume the Stage-1 budget consists of $n_0$ sites
$X_1,\dots,X_{n_0} \sim p_X$, each with $r_0 \ge 2$ replications
$Y_{i,1},\dots,Y_{i,r_0}$. Form the site-wise sample standard
deviation
\[
\tilde\sigma_i
  \;=\; \sqrt{\tfrac{1}{r_0 - 1}\,\textstyle\sum_{j=1}^{r_0}
                    (Y_{i,j} - \bar Y_i)^2},
\qquad
\bar Y_i = \tfrac{1}{r_0}\sum_j Y_{i,j},
\]
and smooth across $X$ via a Nadaraya--Watson estimator with kernel
$K$ and bandwidth $h_\sigma$:
\begin{equation}\label{eq:nw-sigma}
\hat\sigma_{h_\sigma}(x)
  \;=\;
  \frac{\sum_{i=1}^{n_0} K_{h_\sigma}(x - X_i)\,\tilde\sigma_i}
       {\sum_{i=1}^{n_0} K_{h_\sigma}(x - X_i)}.
\end{equation}
Under $Y = f(x) + \sigma(x)\varepsilon$ with
$\mathbb{E}[\varepsilon]=0$, $\mathrm{Var}(\varepsilon)=1$, a
standard delta-method expansion gives
\[
\mathbb{E}[\tilde\sigma_i \mid X_i] = \sigma(X_i) + O(r_0^{-1}),
\quad
\mathrm{Var}(\tilde\sigma_i \mid X_i)
   = \frac{\sigma^2(X_i)\,(\kappa_4 - 1)}{4(r_0-1)} + O(r_0^{-2}),
\]
where $\kappa_4 = \mathbb{E}[\varepsilon^4]$.

\paragraph{Bias--variance decomposition.}
Under the usual NW assumptions (twice-differentiable $\sigma$ and
$p_X$, $p_X(x)>0$, symmetric kernel $K$ with $\mu_2(K)=\int u^2
K(u)\,du$, $R(K) = \int K^2$), the pointwise bias and variance of
$\hat\sigma_{h_\sigma}$ satisfy, as $n_0\to\infty$, $h_\sigma\to
0$, $n_0 h_\sigma \to \infty$:
\begin{equation}\label{eq:sigma-bias}
\boxed{\;
\mathrm{Bias}\bigl[\hat\sigma_{h_\sigma}(x)\bigr]
  \;=\; \tfrac{1}{2}\,h_\sigma^{2}\,\mu_2(K)\!
        \left[\sigma''(x) + 2\,\frac{\sigma'(x)\,p_X'(x)}{p_X(x)}\right]
       + o(h_\sigma^{2}) + O(r_0^{-1}).
\;}
\end{equation}
\begin{equation}\label{eq:sigma-var}
\boxed{\;
\mathrm{Var}\bigl[\hat\sigma_{h_\sigma}(x)\bigr]
  \;=\;
  \frac{R(K)\,\sigma^{2}(x)\,(\kappa_4 - 1)}
       {4\,(r_0-1)\,n_0\,h_\sigma\,p_X(x)}
  \;+\; o\bigl((n_0 h_\sigma r_0)^{-1}\bigr).
\;}
\end{equation}
Two features distinguish~\eqref{eq:sigma-var} from a vanilla NW
variance: (i) the $(r_0-1)^{-1}$ factor reflects the within-site
replication, so doubling $r_0$ halves the variance contribution;
(ii) the kurtosis factor $\kappa_4-1$ makes the variance sensitive
to heavy-tailed noise (e.g. the Gamma or Student-$t$ designs).

\paragraph{Optimal bandwidth and MSE rate.}
The pointwise MSE takes the canonical form
$\mathrm{MSE}(x) = C_1(x)\,h_\sigma^{4}
                    + C_2(x)/(n_0 h_\sigma r_0)$,
and minimising over $h_\sigma$ yields
\begin{equation}\label{eq:hsigma-optimal}
\boxed{\;
h_\sigma^{\star}(x)
   \;\propto\; (n_0\,r_0)^{-1/5},
\qquad
\mathrm{MSE}^{\star}(x)
   \;\propto\; (n_0\,r_0)^{-4/5}.
\;}
\end{equation}
Crucially, the \emph{effective sample size} is $n_0 r_0$, not just
$n_0$: each replicate sharpens $\tilde\sigma_i$ and propagates
through the smoother. Thus under a fixed total budget
$N = n_0 r_0$, the rate of $\hat\sigma$ depends only on $N$, while
Stage-1 CKME fitting depends on $n_0$ alone --- these two
considerations pull $(n_0,r_0)$ in opposite directions.

\paragraph{Propagating $\hat\sigma$ into the coverage bound.}
Plugging~\eqref{eq:nw-sigma} into the adaptive bandwidth rule
$\hat h(x) = c\,\hat\sigma(x)$ and tracking the perturbation
through the score density $g$ in~\eqref{eq:our-error}, a first-order
expansion yields an updated form of the coverage bound of
Proposition~\ref{prop:two-sided}:
\begin{equation}\label{eq:two-sided-plugin}
\bigl|\mathbb{P}(Y\in C(x)\mid X=x) - (1-\alpha)\bigr|
  \;\le\;
  \underbrace{\tfrac{\pi^2 c^{2}}{6}\,\|g''\|_{1}}_{\text{smoothing}}
  \;+\; \underbrace{\|\widehat F - F\|_{\infty}}_{\text{CKME fit}}
  \;+\; \underbrace{c\,\|g'\|_{\infty}\,
        (n_0 r_0)^{-2/5}}_{\text{plug-in $\hat\sigma$}}
  \;+\; p_{n+1}.
\end{equation}
The plug-in term inherits the MSE $^{1/2}$ rate
$(n_0 r_0)^{-2/5}$ from~\eqref{eq:hsigma-optimal} and enters
linearly in $c$: larger $c$ not only amplifies the smoothing bias
but also the plug-in bias.

\paragraph{Implication: the optimal $c$ shrinks under plug-in.}
Let $c_{\text{or}}^{\star}$ minimise the oracle bound
$\tfrac{\pi^2 c^{2}}{6}\|g''\|_{1} + \|\widehat F - F\|_{\infty}$
and $c_{\text{pi}}^{\star}$ minimise~\eqref{eq:two-sided-plugin}.
Ignoring the constant CKME term, differentiating in $c$ gives
\[
c_{\text{pi}}^{\star}
  \;\approx\;
  c_{\text{or}}^{\star}
  \;-\;
  \frac{3\,\|g'\|_{\infty}}{\pi^{2}\,\|g''\|_{1}}
         \,(n_0 r_0)^{-2/5},
\]
so the plug-in optimum is \emph{smaller} than the oracle one, with
a gap that closes at rate $(n_0 r_0)^{-2/5}$. This matches the
empirical observation in the plug-in variant of the homogeneity
experiment, where the best $c$ on the $c$-scan shifts leftwards
relative to the oracle curve, and the gap narrows as we enlarge
$r_0$.

\subsection{Bandwidth Rate for the CKME Indicator}
\label{sec:h-rate}

Proposition~\ref{prop:two-sided} bounds the coverage deviation by
$\tfrac{\pi^{2} c^{2}}{6}\|g''\|_{1} + \|\widehat F - F\|_{\infty}$,
where the first term is a \emph{bias} contribution from
$h(x) = c\,\sigma(x)$ smoothing the score density $g$, while the
second bundles the entire CKME estimation error. Hidden inside
$\|\widehat F - F\|_{\infty}$ is a second, symmetric bias--variance
trade-off in the bandwidth $h$ itself, which we unpack here.

\paragraph{Bias and variance of $\widehat F(y|x)$.}
With a smooth indicator $\psi_{h}(y - t) \approx \mathbf{1}\{y \le t\}$
of bandwidth $h$, the CKME CDF estimator
$\widehat F(y|x) = \sum_i c_i(x)\,\psi_{h}(y - Y_i)$ admits the
pointwise decomposition
\begin{equation}\label{eq:F-bias}
\boxed{\;
\mathbb{E}[\widehat F(y|x)] - F(y|x)
  \;=\;
  \tfrac{1}{2}\,h^{2}\,\mu_{2}(\psi)\,f_{Y|X}'(y|x)
  \;+\; \underbrace{O\!\bigl(\ell_{x}^{2}\bigr)}_{\text{NW bias in }x}
  \;+\; o(h^{2}),
\;}
\end{equation}
\begin{equation}\label{eq:F-var}
\boxed{\;
\mathrm{Var}[\widehat F(y|x)]
  \;=\;
  \frac{F(y|x)(1 - F(y|x))}{n_0\,r_0\,p_X(x)\,\ell_{x}}
  \;+\; O\!\bigl(h^{-1}/(n_0 r_0)\bigr)
  \;+\; o(\cdot).
\;}
\end{equation}
Here $\mu_{2}(\psi) = \int u^{2}\psi'(u)\,du$ plays the role of the
second moment of the induced kernel
$\psi_{h}' = h^{-1}\psi'(\cdot/h)$. Equation~\eqref{eq:F-bias}
says $h$ acts on the CDF exactly as a kernel smoother on $F$, and
its bias is driven by the \emph{score density slope}
$f_{Y|X}'(y|x) = \partial_{y} f_{Y|X}(y|x)$.

\paragraph{Rate for $h^{\star}$.}
Balancing the $h^{4}$ bias² in~\eqref{eq:F-bias} against the
$h^{-1}/(n_0 r_0)$ variance contribution gives the classical
Nadaraya--Watson rate
\begin{equation}\label{eq:h-optimal}
\boxed{\;
h^{\star}\;\propto\;(n_0\,r_0)^{-1/5},
\qquad
\bigl\|\widehat F - F\bigr\|_{\infty}^{\star}\;\propto\;(n_0\,r_0)^{-2/5}.
\;}
\end{equation}
Note the \emph{same} effective sample size $n_0 r_0$ as for
$\hat\sigma$ in~\eqref{eq:hsigma-optimal}: replications help
$\widehat F$ too, because each $Y_{i,j}$ contributes an independent
vote to the indicator average. This is why the replicated two-stage
design $(n_0, r_0)$ with $r_0 \ge 2$ dominates a single-shot design
of the same total $N = n_0 r_0$ in the tail-quantile regime (where
variance dominates).

\paragraph{Combining with the coverage bound.}
Substituting~\eqref{eq:h-optimal} into the coverage decomposition
and collecting terms, Proposition~\ref{prop:two-sided} sharpens to
\begin{equation}\label{eq:two-sided-full}
\bigl|\mathbb{P}(Y\in C(x)\mid X=x) - (1-\alpha)\bigr|
  \;\lesssim\;
  \underbrace{\tfrac{\pi^{2} c^{2}}{6}\|g''\|_{1}}_{\text{$c$-smoothing}}
  \;+\; \underbrace{(n_0 r_0)^{-2/5}}_{\text{$h$ trade-off}}
  \;+\; \underbrace{c\,(n_0 r_0)^{-2/5}\|g'\|_{\infty}}_{\text{$\hat\sigma$ plug-in}}
  \;+\; p_{n+1}.
\end{equation}
Three distinct bias--variance mechanisms are visible:
(i) the homogeneity term driven by $c$, which \emph{cannot} be
driven to zero at finite $n$ because $c\to 0$ sends the adaptive
mechanism off; (ii) the CKME bandwidth $h$ term, which vanishes at
rate $(n_0 r_0)^{-2/5}$; and (iii) the plug-in $\hat\sigma$ term,
also at $(n_0 r_0)^{-2/5}$ but scaled by $c$.

\begin{remark}[Why high-$\tau$ quantiles do not converge at fixed $h$]
\label{rem:tail-nonconvergence}
Inverting~\eqref{eq:F-bias} gives
$\hat q_{\tau}(x) - q_{\tau}(x) \approx
    -\tfrac{1}{2} h^{2}\mu_{2}(\psi)\,
        f_{Y|X}'(q_{\tau}|x) / f_{Y|X}(q_{\tau}|x)$,
so the quantile bias is amplified by the reciprocal tail density
$1/f_{Y|X}(q_{\tau}|x)$. When $\tau\to 1$ this factor diverges for
any distribution with a density tail, so \emph{fixing} $h$ (no
matter how small) leaves a floor in $\sup_{x}|\hat q_{\tau}-q_{\tau}|$
that does not vanish as $n\to\infty$. Only $h\to 0$ at the rate
prescribed by~\eqref{eq:h-optimal} restores convergence, and even
then only if $h^{2}/f_{Y|X}(q_{\tau}|x)\to 0$ ---
a $\tau$-dependent condition. This matches the empirical finding
in \texttt{exp\_onesided/exp2\_sup\_vs\_tau.py}, where
$\sup_{x}|\hat q_{\tau}-q_{\tau}|$ decays with $n$ for
$\tau\in\{0.5,0.6,0.7\}$ but stalls for
$\tau\in\{0.8,0.9,0.95\}$ under the project's default fixed $h$.
\end{remark}

\paragraph{Summary: three U-shapes, one budget.}
Equation~\eqref{eq:two-sided-full} predicts \emph{three}
bias--variance U-shapes along the CKME pipeline: one in the CKME
bandwidth $h$, one in the adaptive coefficient $c$, and one in the
plug-in smoother bandwidth $h_{\sigma}$. All three are confirmed
empirically in Figure~\ref{fig:three-u-shapes}: panel (a) varies
$h$ at fixed $(c,\hat\sigma)$ and shows quantile error as a
function of $h$ across multiple $\tau$ levels; panel (b) varies
$c$ in $h(x)=c\,\hat\sigma(x)$ and reports the resulting
quantile and coverage errors; panel (c) varies the plug-in
smoother bandwidth $h_{\sigma}$ and reports RMSE of $\hat\sigma$,
KS-max of the score, and the sup coverage gap. All three share
the effective sample size $n_0 r_0$, so budget allocation ---
not just parameter tuning --- becomes the correct lens: the
marginal return of one more replication equals that of one more
site whenever variance dominates at that position on the U-curve.

\begin{figure}[!ht]
\centering
\includegraphics[width=0.62\textwidth]{../exp_h_sweep/output_exp2/fig3_qerror_vs_h.png}\\[2pt]
{\small (a) CKME indicator bandwidth $h$ (exp2): quantile error
\(|\hat q_{\tau}-q_{\tau}|\) vs.\ $h$ across $\tau$.}\\[3pt]
\includegraphics[width=0.62\textwidth]{../exp_conditional_coverage/diagnostics/output/diag3_cscan.png}\\[2pt]
{\small (b) Adaptive coefficient $c$: quantile and coverage errors
vs.\ $c$ in $h(x)=c\,\hat\sigma(x)$.}\\[3pt]
\includegraphics[width=0.62\textwidth]{../exp_score_homogeneity_plugin/output/figD_hsigma_scan.png}\\[2pt]
{\small (c) Plug-in smoother bandwidth $h_{\sigma}$: RMSE of
$\hat\sigma$, KS-max, and sup coverage gap vs.\ $h_{\sigma}$ for
three simulators.}
\caption{\textbf{Three bias--variance U-shapes along the CKME
pipeline.} Each panel isolates one bandwidth while holding the
others fixed. The left branch of every U is variance-dominated and
the right branch is bias-dominated, consistent with
Equations~\eqref{eq:F-bias}--\eqref{eq:F-var},
\eqref{eq:two-sided-full} and \eqref{eq:sigma-bias}--\eqref{eq:sigma-var}.}
\label{fig:three-u-shapes}
\end{figure}

% ------------------------------------------------------------------
\subsection{Generalization Beyond Gaussian Noise}
\label{sec:beyond-gaussian}
% ------------------------------------------------------------------

The effective-bandwidth argument (Section~\ref{sec:h-eff}) uses only the
location-scale structure $Y = f(x) + \sigma(x)\varepsilon$, not the
specific form of $G$. The key algebraic step
\[
  \frac{y - Y_{ij}}{h(x)} = \frac{z - Z_{ij}}{c}
\]
holds regardless of whether $\varepsilon$ is Gaussian, Student-$t$,
centered Gamma, or any other distribution with $E[\varepsilon] = 0$,
$\mathrm{Var}(\varepsilon) = 1$.

Therefore:
\begin{itemize}
  \item Proposition~\ref{prop:homogeneity} applies to all
    location-scale families.
  \item The proof requires only that $G$ has a bounded density (for
    kernel smoothing consistency), not that $G$ is Gaussian.
  \item Beyond location-scale (e.g., shape parameter varying with $x$),
    the standardization argument breaks. However, experiments show
    approximate homogeneity is preserved for moderate violations
    (Section~\ref{sec:exp-design}).
\end{itemize}


%% ====================================================================
\section{Empirical Validation}
\label{sec:exp-design}
%% ====================================================================

\subsection{Objective}

The experiment tests three theoretical predictions:
\begin{enumerate}
  \item Adaptive $h(x) = c\,\sigma(x)$ produces lower KS distance
    between local and pooled score distributions (score homogeneity).
  \item Score homogeneity, not estimation accuracy, determines
    conditional coverage gap (decoupling).
  \item The mechanism holds for non-Gaussian location-scale families
    (generality).
\end{enumerate}

\subsection{Three-Phase Design}

We use three simulators with increasing difficulty, all sharing the
true function $f(x) = e^{x/10}\sin(x)$ on $x \in [0, 2\pi]$:

\begin{center}
\begin{tabular}{@{}clll@{}}
\toprule
\textbf{Phase} & \textbf{Simulator} & \textbf{$\sigma(x)$ profile} &
  \textbf{Purpose} \\
\midrule
1 & \texttt{mild} &
  $0.3(1 + 0.5\sin x)$, ratio $\approx 3\times$ &
  Clean theory verification \\
2 & \texttt{exp2} &
  $\sqrt{0.01 + 0.2(x-\pi)^2}$, ratio $\approx 200\times$ &
  Finite-sample stress test \\
3 & \texttt{nongauss\_B2L} &
  $\sigma_{\mathrm{tar}}(x) = 0.1 + 0.1(x-\pi)^2$, Gamma($k$=2) &
  Non-Gaussian robustness \\
\bottomrule
\end{tabular}
\end{center}

\paragraph{Why three phases?}
Phase~1 (mild) has a small $\sigma$-ratio, so finite-sample noise is
small and the effective-bandwidth mechanism can be verified cleanly.
Phase~2 (exp2) has a $\sigma$-ratio of $\sim$200$\times$, stressing the
mechanism where $\sigma(x) \to 0$ creates near-degenerate CDF
estimates. Phase~3 (nongauss\_B2L) uses centered Gamma noise ($k=2$,
strong skew), testing whether the location-scale argument extends to
non-Gaussian settings.

\subsection{Bandwidth Configurations}

Nine bandwidth configurations per simulator:
\begin{itemize}
  \item \textbf{3 fixed:} $h = 0.05$ (small), $h_{\mathrm{cv}}$
    (cross-validated), $h = 0.5$ (large).
  \item \textbf{6 adaptive:} $h(x) = c\,\sigma(x)$ with oracle
    $\sigma(x)$, $c \in \{0.5, 1.0, 1.5, 2.0, 3.0, 5.0\}$.
\end{itemize}

All configurations share the same $(e\kern-0.5pt\ell_x, \lambda)$
hyperparameters (cross-validated jointly), isolating the bandwidth
effect.

\subsection{Budget and Evaluation}

\begin{itemize}
  \item \textbf{Stage~1:} $n_0 = 512$ sites, $r_0 = 1$ replication
    each. Stage~2: $n_1 = 512$, $r_1 = 1$.
  \item \textbf{Evaluation grid:} $M = 50$ equally-spaced points in
    $[0, 2\pi]$.
  \item \textbf{MC test draws:} $B = 2000$ fresh draws per eval point
    per macrorep for conditional coverage and score ECDF.
  \item \textbf{Macroreplications:} $n_{\mathrm{macro}} = 20$
    (5 completed for current results).
\end{itemize}

\subsection{Metrics}

\begin{enumerate}
  \item \textbf{KS$_{\max}$}: $\max_{m=1}^{M}
    \mathrm{KS}(F_{s|x_m}, F_{s,\mathrm{pooled}})$ --- score
    homogeneity (lower = better).
  \item \textbf{KS$_{\mathrm{mean}}$}: $\frac{1}{M}\sum_m
    \mathrm{KS}(F_{s|x_m}, F_{s,\mathrm{pooled}})$ --- average
    homogeneity.
  \item \textbf{Coverage gap (sup)}:
    $\max_m |P(Y \in C(x_m)|x_m) - (1-\alpha)|$ --- worst-case
    conditional coverage deviation (lower = better).
  \item \textbf{Coverage gap (MAE)}:
    $\frac{1}{M}\sum_m |P(Y \in C(x_m)|x_m) - (1-\alpha)|$ ---
    average coverage deviation.
  \item \textbf{Quantile error (sup)}:
    $\max_m \max_{\tau \in \{\alpha/2, 1-\alpha/2\}}
    |\hat{q}_\tau(x_m) - q_\tau(x_m)|$ --- estimation accuracy
    (not directly related to coverage under homogeneity).
  \item \textbf{Width std}: $\mathrm{SD}_x(W(x_m))$ --- spatial
    variation of interval width (higher = more locally adaptive).
\end{enumerate}


%% ====================================================================
% NOTE-ONLY: detailed per-phase results; in paper mode, condense to key figures
\ifpaper\else
\section{Detailed Experiment Results}
\label{sec:results}
%% ====================================================================

All results are mean $\pm$ SD across 5 macroreplications.

\newcommand{\figdir}{../exp_score_homogeneity/output}

% ------------------------------------------------------------------
\subsection{Phase 1: Mild Heteroscedasticity}
\label{sec:results-mild}
% ------------------------------------------------------------------

\begin{table}[h]
\centering
\small
\caption{Phase~1 (\texttt{mild}): selected bandwidth configurations.}
\label{tab:mild}
\begin{tabular}{@{}lccccc@{}}
\toprule
\textbf{Bandwidth} & \textbf{KS$_{\max}$} & \textbf{KS$_{\mathrm{mean}}$}
  & \textbf{Cov gap (sup)} & \textbf{Cov gap (MAE)} & \textbf{$q$-err (sup)} \\
\midrule
fixed\_small & $0.236 \pm 0.028$ & $0.089 \pm 0.014$ & $0.105 \pm 0.042$
  & $0.039 \pm 0.005$ & $1.83 \pm 0.08$ \\
fixed\_cv    & $0.189 \pm 0.004$ & $0.114 \pm 0.010$ & $0.115 \pm 0.028$
  & $0.064 \pm 0.010$ & $1.83 \pm 0.08$ \\
fixed\_large & $0.267 \pm 0.007$ & $0.147 \pm 0.003$ & $0.131 \pm 0.011$
  & $0.075 \pm 0.005$ & $1.83 \pm 0.08$ \\
\midrule
adaptive $c\!=\!0.5$ & $0.142 \pm 0.034$ & $0.050 \pm 0.008$ & $0.063 \pm 0.008$
  & $0.023 \pm 0.003$ & $1.83 \pm 0.08$ \\
adaptive $c\!=\!1.0$ & $0.081 \pm 0.011$ & $0.031 \pm 0.004$ & $0.047 \pm 0.008$
  & $0.013 \pm 0.002$ & $1.83 \pm 0.08$ \\
adaptive $c\!=\!1.5$ & $\mathbf{0.065 \pm 0.019}$ & $\mathbf{0.025 \pm 0.002}$
  & $\mathbf{0.042 \pm 0.013}$ & $\mathbf{0.010 \pm 0.002}$
  & $1.83 \pm 0.08$ \\
adaptive $c\!=\!2.0$ & $0.069 \pm 0.029$ & $0.023 \pm 0.001$ & $0.047 \pm 0.018$
  & $0.010 \pm 0.003$ & $2.04 \pm 0.06$ \\
adaptive $c\!=\!5.0$ & $0.160 \pm 0.086$ & $0.026 \pm 0.003$ & $0.121 \pm 0.071$
  & $0.013 \pm 0.003$ & $2.93 \pm 0.03$ \\
\bottomrule
\end{tabular}
\end{table}

\paragraph{Finding~1: Adaptive $h$ dramatically improves score homogeneity.}
The best adaptive configuration ($c = 1.5$) achieves
$\mathrm{KS}_{\max} = 0.065$, compared to the best fixed
($\mathrm{fixed\_cv}$, $0.189$) --- a $3\times$ reduction. This
directly validates Proposition~\ref{prop:homogeneity}.

\paragraph{Finding~2: Coverage gap tracks homogeneity, not estimation
accuracy.}
All fixed bandwidths have similar $q$-error (sup) $\approx 1.83$ (they
share the same $e\kern-0.5pt\ell_x, \lambda$), yet
$\mathrm{fixed\_large}$ has the worst coverage gap ($0.131$) while
$\mathrm{fixed\_small}$ has a smaller gap ($0.105$). Meanwhile,
adaptive $c = 1.5$ has the same $q$-error but the lowest coverage gap
($0.042$). This confirms that \emph{homogeneity, not estimation
accuracy, determines conditional coverage}.

\begin{figure}[ht]
\centering
\includegraphics[width=0.85\textwidth]{\figdir/fig1_score_homogeneity_profile_mild.png}
\caption{Phase~1 (\texttt{mild}): KS distance from pooled score
  distribution vs.\ $x$, comparing \texttt{fixed\_cv} (left) and
  \texttt{adaptive\_c2.0} (right). Mean $\pm$ 1SD across 5 macroreps.
  Adaptive bandwidth produces a flatter KS profile, indicating more
  homogeneous scores across $x$.}
\label{fig:homog-mild}
\end{figure}

\begin{figure}[ht]
\centering
\includegraphics[width=0.85\textwidth]{\figdir/fig2_coverage_curve_mild.png}
\caption{Phase~1 (\texttt{mild}): conditional coverage vs.\ $x$ for
  selected bandwidth configurations. Mean $\pm$ 1SD across 5 macroreps.
  Dashed line = target $1 - \alpha = 0.9$. Adaptive bandwidths maintain
  coverage closer to the target across the full input range.}
\label{fig:coverage-mild}
\end{figure}

% ------------------------------------------------------------------
\subsection{Phase 2: Extreme Heteroscedasticity (\texttt{exp2})}
\label{sec:results-exp2}
% ------------------------------------------------------------------

\begin{table}[h]
\centering
\small
\caption{Phase~2 (\texttt{exp2}): selected bandwidth configurations.}
\label{tab:exp2}
\begin{tabular}{@{}lccccc@{}}
\toprule
\textbf{Bandwidth} & \textbf{KS$_{\max}$} & \textbf{KS$_{\mathrm{mean}}$}
  & \textbf{Cov gap (sup)} & \textbf{Cov gap (MAE)} & \textbf{$q$-err (sup)} \\
\midrule
fixed\_small & $0.806 \pm 0.006$ & $0.264 \pm 0.008$ & $0.171 \pm 0.058$
  & $0.065 \pm 0.010$ & $1.36 \pm 0.58$ \\
fixed\_cv    & $0.806 \pm 0.003$ & $0.289 \pm 0.006$ & $0.168 \pm 0.057$
  & $0.076 \pm 0.011$ & $1.33 \pm 0.57$ \\
fixed\_large & $0.808 \pm 0.013$ & $0.341 \pm 0.005$ & $0.245 \pm 0.062$
  & $0.099 \pm 0.008$ & $1.62 \pm 0.22$ \\
\midrule
adaptive $c\!=\!1.0$ & $0.743 \pm 0.026$ & $0.182 \pm 0.003$ & $0.114 \pm 0.030$
  & $0.048 \pm 0.015$ & $2.48 \pm 0.08$ \\
adaptive $c\!=\!2.0$ & $0.661 \pm 0.025$ & $0.136 \pm 0.003$ & $0.113 \pm 0.029$
  & $0.037 \pm 0.013$ & $3.86 \pm 0.17$ \\
adaptive $c\!=\!5.0$ & $\mathbf{0.479 \pm 0.004}$ & $\mathbf{0.076 \pm 0.003}$
  & $0.153 \pm 0.093$ & $\mathbf{0.031 \pm 0.010}$ & $4.78 \pm 0.26$ \\
\bottomrule
\end{tabular}
\end{table}

\paragraph{Finding~3: Extreme $\sigma$-ratio reveals finite-sample limits.}
With $\sigma$-ratio $\approx 200\times$, even the best adaptive
configuration ($c = 5.0$) achieves $\mathrm{KS}_{\max} = 0.479$,
much higher than Phase~1. All fixed bandwidths have
$\mathrm{KS}_{\max} > 0.8$, reflecting near-total score heterogeneity.
The near-zero $\sigma(\pi) \approx 0.1$ creates degenerate CDF
estimates that limit homogeneity improvement.

\paragraph{Finding~4: Decoupling is even more pronounced.}
The $q$-error (sup) \emph{increases} from $c = 0.5$ ($1.23$) to
$c = 5.0$ ($4.78$), while the coverage MAE \emph{decreases} from
$0.052$ to $0.031$. This anti-diagonal pattern between estimation
accuracy and conditional coverage is the hallmark of the decoupling
predicted by the score homogeneity theory.

\begin{figure}[ht]
\centering
\includegraphics[width=0.85\textwidth]{\figdir/fig1_score_homogeneity_profile_exp2.png}
\caption{Phase~2 (\texttt{exp2}): KS distance from pooled score
  distribution vs.\ $x$. The extreme $\sigma$-ratio ($\sim$200$\times$)
  causes high KS values even for adaptive bandwidth, especially near
  $x = \pi$ where $\sigma(x) \to 0$.}
\label{fig:homog-exp2}
\end{figure}

\begin{figure}[ht]
\centering
\includegraphics[width=0.85\textwidth]{\figdir/fig2_coverage_curve_exp2.png}
\caption{Phase~2 (\texttt{exp2}): conditional coverage vs.\ $x$.
  Fixed bandwidths show large coverage deviations; adaptive bandwidths
  reduce the gap but cannot fully resolve the near-zero-$\sigma$ region
  at $x \approx \pi$.}
\label{fig:coverage-exp2}
\end{figure}

% ------------------------------------------------------------------
\subsection{Phase 3: Non-Gaussian (\texttt{nongauss\_B2L})}
\label{sec:results-B2L}
% ------------------------------------------------------------------

\begin{table}[h]
\centering
\small
\caption{Phase~3 (\texttt{nongauss\_B2L}): selected bandwidth configurations.}
\label{tab:B2L}
\begin{tabular}{@{}lccccc@{}}
\toprule
\textbf{Bandwidth} & \textbf{KS$_{\max}$} & \textbf{KS$_{\mathrm{mean}}$}
  & \textbf{Cov gap (sup)} & \textbf{Cov gap (MAE)} & \textbf{$q$-err (sup)} \\
\midrule
fixed\_small & $0.459 \pm 0.022$ & $0.182 \pm 0.007$ & $0.156 \pm 0.021$
  & $0.060 \pm 0.005$ & $1.04 \pm 0.61$ \\
fixed\_cv    & $0.445 \pm 0.017$ & $0.189 \pm 0.006$ & $0.177 \pm 0.020$
  & $0.064 \pm 0.004$ & $1.02 \pm 0.61$ \\
fixed\_large & $0.485 \pm 0.013$ & $0.278 \pm 0.003$ & $0.287 \pm 0.053$
  & $0.098 \pm 0.006$ & $1.46 \pm 0.12$ \\
\midrule
adaptive $c\!=\!1.0$ & $0.263 \pm 0.032$ & $0.082 \pm 0.003$ & $0.072 \pm 0.010$
  & $0.028 \pm 0.005$ & $1.81 \pm 0.10$ \\
adaptive $c\!=\!1.5$ & $\mathbf{0.172 \pm 0.024}$ & $0.058 \pm 0.005$
  & $\mathbf{0.070 \pm 0.025}$ & $\mathbf{0.023 \pm 0.004}$
  & $2.34 \pm 0.04$ \\
adaptive $c\!=\!2.0$ & $0.167 \pm 0.023$ & $\mathbf{0.048 \pm 0.006}$
  & $0.095 \pm 0.034$ & $0.021 \pm 0.006$ & $2.89 \pm 0.21$ \\
adaptive $c\!=\!5.0$ & $0.328 \pm 0.037$ & $0.051 \pm 0.006$
  & $0.293 \pm 0.033$ & $0.025 \pm 0.009$ & $3.82 \pm 0.35$ \\
\bottomrule
\end{tabular}
\end{table}

\paragraph{Finding~5: Score homogeneity extends to non-Gaussian noise.}
Despite Gamma($k$=2) noise (strong positive skew), adaptive $c = 1.5$
achieves $\mathrm{KS}_{\max} = 0.172$, compared to the best fixed
($0.445$) --- a $2.6\times$ reduction. This validates the
location-scale generalization of Section~\ref{sec:beyond-gaussian}.

\paragraph{Finding~6: Optimal $c$ shifts slightly for non-Gaussian noise.}
The optimal $c$ for KS$_{\max}$ is $c \approx 1.5$--$2.0$ (non-Gaussian),
compared to $c \approx 1.0$--$1.5$ (Gaussian, Phase~1). The skewed
noise CDF requires slightly more smoothing for the kernel estimate
$\hat{G}_c(z)$ to be stable.

\begin{figure}[ht]
\centering
\includegraphics[width=0.85\textwidth]{\figdir/fig2_coverage_curve_nongauss_B2L.png}
\caption{Phase~3 (\texttt{nongauss\_B2L}): conditional coverage vs.\
  $x$. Despite non-Gaussian (Gamma) noise, adaptive bandwidths achieve
  near-uniform coverage, confirming the location-scale generalization.}
\label{fig:coverage-B2L}
\end{figure}

% ------------------------------------------------------------------
\subsection{Cross-Phase: Pareto and Decoupling Plots}
\label{sec:results-pareto}
% ------------------------------------------------------------------

\begin{figure}[ht]
\centering
\includegraphics[width=0.48\textwidth]{\figdir/fig3_cscan_pareto_mild.png}
\hfill
\includegraphics[width=0.48\textwidth]{\figdir/fig3_cscan_pareto_nongauss_B2L.png}
\caption{Pareto plot: KS$_{\max}$ (score homogeneity) vs.\ coverage gap
  (sup) for Phase~1 (\texttt{mild}, left) and Phase~3
  (\texttt{nongauss\_B2L}, right). Circles = adaptive, squares = fixed.
  Adaptive configurations dominate the lower-left (better) corner.
  Mean $\pm$ 1SD error bars across 5 macroreps.}
\label{fig:pareto}
\end{figure}

\begin{figure}[ht]
\centering
\includegraphics[width=0.48\textwidth]{\figdir/fig4_decoupling_mild.png}
\hfill
\includegraphics[width=0.48\textwidth]{\figdir/fig4_decoupling_nongauss_B2L.png}
\caption{Decoupling plot: quantile error (sup) vs.\ coverage gap (sup)
  for Phase~1 (left) and Phase~3 (right). The anti-diagonal pattern
  confirms decoupling: adaptive bandwidths trade estimation accuracy
  for better conditional coverage through score homogeneity.}
\label{fig:decoupling}
\end{figure}
\fi % end \ifpaper\else for Detailed Experiment Results


%% ====================================================================
\section{Discussion}
\label{sec:discussion}
%% ====================================================================

\subsection{Theory--Experiment Alignment}

The experiments provide consistent support for all three theoretical
predictions (based on $n_{\mathrm{macro}} = 5$; further macroreplications
are planned):
\begin{enumerate}
  \item \textbf{Proposition~\ref{prop:homogeneity}:}
    Adaptive $h$ reduces KS$_{\max}$ by $2.6\times$--$3\times$
    vs.\ best fixed (Phase~1 and~3). Phase~2's smaller improvement
    ($1.7\times$) reflects the finite-sample limitation where
    $\sigma(x) \to 0$.
  \item \textbf{Decoupling (Section~\ref{sec:two-conditions}):}
    The anti-diagonal pattern in the decoupling plots ---
    estimation accuracy vs.\ coverage gap --- is consistently observed
    across all three phases. Fixed bandwidths cluster in the
    ``accurate but inhomogeneous'' corner; adaptive bandwidths move
    toward the ``less accurate but homogeneous'' corner with better
    coverage.
  \item \textbf{Location-scale generality (Section~\ref{sec:beyond-gaussian}):}
    Phase~3 (Gamma noise) shows the same qualitative pattern as
    Phase~1 (Gaussian noise), confirming that the mechanism depends
    on the location-scale structure, not the specific noise
    distribution.
\end{enumerate}

\subsection{Finite-Sample Limitations}

Phase~2 (\texttt{exp2}) reveals important limits:
\begin{itemize}
  \item When $\sigma(x) \to 0$ (at $x \approx \pi$), even adaptive
    $h(x) = c\,\sigma(x) \to 0$ creates a near-degenerate CDF with
    $O(h(x))$ resolution, and the KS distance from the pooled
    distribution remains high.
  \item Very large $c$ ($c = 5.0$) achieves the best KS$_{\max}$
    but the worst coverage gap (sup), illustrating the transition
    toward the over-smoothed collapse regime
    (Proposition~\ref{prop:oversmooth}).
  \item The $c$-scan Pareto plot reveals an optimal range
    $c \approx 1$--$2$ that balances homogeneity and local adaptivity.
\end{itemize}

\subsection{Practical Implications}

\begin{enumerate}
  \item \textbf{Bandwidth selection:} Based on the current three-phase
    evidence, $c \approx 1.5$ appears to be a reasonable default.
    The coverage gap is relatively insensitive to $c$ in the range
    $[1, 2]$ across all three phases.
  \item \textbf{Plug-in estimation:} In practice, $\sigma(x)$ is
    unknown. The two-stage design provides replication-based estimates
    $\hat\sigma(x_j)$ from Stage~1, enabling
    $\hat{h}(x) = c\,\hat\sigma(x)$. The simulation evidence
    motivates this plug-in approach, but quantifying how
    $\hat\sigma(x)$ estimation error degrades the homogeneity bound
    is addressed by the rate analysis in
    Section~\ref{sec:plugin-rate}.
  \item \textbf{CDF-first advantage:} The score homogeneity mechanism
    is specific to CDF-based scores (Remark~\ref{rem:cdf-first}).
    Quantile-first methods (CQR, DCP-QR) cannot exploit this regime,
    even with adaptive bandwidth.
\end{enumerate}

\subsection{Comparison with Related Theoretical Frameworks}
\label{sec:related-comparison}

We position our score homogeneity results relative to two recent
frameworks for conditional conformal prediction: CP$^2$ \cite{CP2}
and RCP \cite{RCP}. Table~\ref{tab:comparison} summarizes the
comparison.

\paragraph{vs.\ CP$^2$ (Gibbs, Cherian, Cand\`{e}s 2023).}
CP$^2$ provides a general conditional coverage framework via RKHS
regularization. It bounds the conditional coverage gap by
$\|f\|_{\mathcal{H}} \cdot R_n$, where $\|f\|_{\mathcal{H}}$ is a
score distribution regularity term and $R_n$ a calibration error.
This is a powerful ``black-box'' result: it applies broadly but does
not identify \emph{when} the score distribution is well-behaved, nor
does it connect to any specific method parameter.

Our contribution fills this gap in two ways:
\begin{enumerate}
  \item \textbf{Sufficient condition.} We identify a concrete,
    verifiable condition --- $h(x) = c\,\sigma(x)$ under a
    location-scale model --- that makes the CDF-based score
    distribution $x$-independent.
  \item \textbf{Explicit quantitative bound.} Using the classical
    result of Devroye and Gy\"{o}rfi (1985) \cite{DG85} for the
    logistic kernel, we obtain
    \[
      d_{\mathrm{TV}}(\psi_c * g,\; g)
        \le \frac{\pi^2 c^2}{6}\,\|g''\|_1,
    \]
    where $\psi_c$ is the logistic kernel scaled by $c$. This yields
    $d_{\mathrm{TV}}(F_{s|x}, F_{s,\mathrm{pivot}})
      \le \frac{\pi^2 c^2}{6}\,\|g''\|_1 + O_p(n^{-s})$,
    separating smoothing bias from regression error.
    The bound is computable: given $c$ and the noise density $g$, one
    evaluates $\|g''\|_1$ directly.
\end{enumerate}

In short, CP$^2$ says ``if the score distribution is regular enough,
coverage is good''; we say ``under adaptive bandwidth, the score
distribution \emph{is} regular, and here is the explicit bound.''

\paragraph{vs.\ RCP (Rectified Conformal Prediction).}
RCP addresses the same problem --- score heterogeneity degrades
conditional coverage --- but operates on a fundamentally different
error decomposition (see Section~\ref{sec:error-comparison} for the
formal comparison). RCP's coverage bound is controlled by the
\emph{quantile-of-score estimation error}
$\sup_x |\hat\tau(x) - \tau(x)|$, where $\hat\tau(x)$ estimates the
conditional quantile of the rectified score. Our coverage is instead
controlled by $d_{\mathrm{TV}}(F_{s|x}, F_{s,\mathrm{pivot}})$, which
decomposes into smoothing bias $O(c^2)$ and CDF estimation error
$O_p(n^{-s})$.

The practical consequence is that our approach achieves score
homogeneity \emph{without any post-hoc transformation}: under
adaptive bandwidth, the CDF-first scores are structurally
$x$-independent (``natural rectification''), whereas RCP requires
(i)~estimating a rectification map $T$, (ii)~additional sample
splitting, and (iii)~Lipschitz regularity conditions on $T$.
This is specific to CDF-based scores --- quantile-first methods
(CQR, DCP-QR) do not admit this mechanism.

\paragraph{Additional elements.}
Two aspects of our analysis, to our knowledge, are not addressed by
either CP$^2$ or RCP:

\begin{enumerate}
  \item \textbf{Two-condition framework (Section~\ref{sec:two-conditions}).}
    To our knowledge, neither CP$^2$ nor RCP explicitly distinguishes
    between ``useful'' and ``trivial'' score homogeneity. We show that
    $h \to \infty$ also
    produces $x$-independent scores (Proposition~\ref{prop:oversmooth}),
    but with degenerate constant-width intervals. Our joint requirement
    of homogeneity \emph{plus} local adaptivity
    ($W(x)/\sigma(x) \to \mathrm{const}$) excludes this trivial
    regime and pins down the useful range $c \in [1, 2]$.
  \item \textbf{Explicit bias--variance decomposition in $c$.}
    The total $d_{\mathrm{TV}}$ decomposes as
    $\underbrace{\tfrac{\pi^2 c^2}{6}\|g''\|_1}_{\text{smoothing bias}}
    + \underbrace{O_p(n^{-s})}_{\text{regression error}}$.
    Small $c$ reduces smoothing bias but increases regression error
    (CDF estimate variance); large $c$ does the reverse. This
    decomposition provides an explicit, operational guide for bandwidth
    selection in the conformal prediction setting, and it is consistent
    with the experimentally observed optimal range.
\end{enumerate}

\begin{table}[htbp]
\centering
\caption{Comparison of conditional coverage frameworks.}
\label{tab:comparison}
\small
\begin{tabular}{@{}lccc@{}}
\toprule
  & CP$^2$ & RCP & Ours \\
\midrule
Sufficient condition for homogeneity
  & \texttimes & via $T(s,x)$ & $h = c\,\sigma(x)$ \\
Quantitative bound
  & implicit (RKHS) & Lipschitz & $\pi^2 c^2/6 \cdot \|g''\|_1$ \\
Extra estimation step
  & --- & estimate $T$ & none \\
Excludes over-smoothed collapse
  & \texttimes & \texttimes & \checkmark \\
Bandwidth selection guidance
  & \texttimes & \texttimes & \checkmark \\
Scope
  & general & general & location-scale \\
\bottomrule
\end{tabular}
\end{table}

\paragraph{Scope trade-off.}
Our results apply to the location-scale family
$Y = f(x) + \sigma(x)\varepsilon$, which is narrower than the
settings of CP$^2$ or RCP. Within this class, however, the analysis yields more explicit and
operational conclusions: a verifiable sufficient condition, a
computable bound, and a principled bandwidth selection rule. For
applications where location-scale is a reasonable model (simulation,
engineering reliability, heteroscedastic regression), these concrete
results complement the broader but more implicit guarantees of
CP$^2$ and RCP.


%% ====================================================================
% NOTE-ONLY: development memo, not for paper
\ifpaper\else
\section{Connection to Full Paper}
\label{sec:connection}
%% ====================================================================

This note expands Section~4 of \texttt{paper\_strategy\_memo.tex}
(Result~2: Score Homogeneity Theorem). For the paper version:

\begin{itemize}
  \item \textbf{Keep:} Theorem statement, effective bandwidth argument,
    over-smoothed collapse proposition, CP$^2$ connection, and summary
    results table (one selected table, e.g., Phase~1).
  \item \textbf{Condense:} The three-phase experiment can be summarized
    in one paragraph with a reference to the supplementary material.
    The full tables (Tables~\ref{tab:mild}--\ref{tab:B2L}) move to an
    appendix.
  \item \textbf{Add for paper:} Formal proof of
    Theorem~\ref{prop:homogeneity} (currently proof sketch), plug-in
    stability proposition, and comparison with RCP's rectification
    approach.
\end{itemize}

\paragraph{Open questions.}
\begin{enumerate}
  \item Can the optimal $c$ be determined from the noise distribution
    $G$ without a grid search?
  \item Does the homogeneity mechanism extend to higher-dimensional
    inputs ($d > 1$)?
  \item How does plug-in $\hat\sigma(x)$ estimation error degrade the
    homogeneity bound quantitatively?
\end{enumerate}
\fi % end \ifpaper\else for Connection to Full Paper


%% ====================================================================
%% References (placeholder)
%% ====================================================================
\begin{thebibliography}{9}
\bibitem{CP2}
  Gibbs, I., Cherian, J.J., and Cand\`{e}s, E.J. (2023).
  Conformal prediction with conditional guarantees.
  \emph{arXiv:2305.12616}.

\bibitem{RCP}
  Bhatnagar, A., Wang, H., Xiong, C., and Bai, Y. (2025).
  Improved conditional coverage via rectified conformity scores.
  \emph{arXiv:2502.16336v2}.

\bibitem{Barber2023}
  Barber, R.F., Cand\`{e}s, E.J., Ramdas, A., and Tibshirani, R.J. (2023).
  Conformal prediction beyond exchangeability.
  \emph{Annals of Statistics}, 51(2), 816--845.

\bibitem{DG85}
  Devroye, L. and Gy\"{o}rfi, L. (1985).
  \emph{Nonparametric Density Estimation: The $L_1$ View}.
  Wiley, New York.
\end{thebibliography}

\end{document}
